\section{Bố cục báo cáo}
Nội dung còn lại của báo cáo này được tổ chức thành các chương như sau:

\begin{itemize}
    \item \textbf{Chương 2: Cơ sở lý thuyết (Theoretical Background):} Trình bày các nền tảng toán học và kỹ thuật cơ bản, bao gồm Structure-from-Motion (SfM), trường độ mờ Gaussian (Gaussian Opacity Fields - GOF), các thuật toán trích xuất lưới tam giác, và nguyên lý hoạt động của các engine mô phỏng vật lý.
    
    \item \textbf{Chương 3: Tập dữ liệu và các độ đo (Datasets and Metrics):} Trình bày chi tiết các bộ dữ liệu dùng để thử nghiệm — bao gồm tập dữ liệu tổng hợp trong nhà \textit{Replica}, tập dữ liệu gắn nhãn ngôn ngữ \textit{LERF}, và cảnh chụp thực tế \textit{teatime} — cùng hệ thống các độ đo về chất lượng ảnh (PSNR, SSIM), chất lượng hình học (Chamfer Distance) và độ ổn định va chạm vật lý.
    
    \item \textbf{Chương 4: Các công trình liên quan (Related Works):} Điểm qua các nghiên cứu tiêu biểu về 3D Gaussian Splatting, các phương pháp chỉnh sửa và xóa vật thể 3D, kỹ thuật hoàn thiện 3D khi bị che khuất (Amodal 3D Completion), và ứng dụng của mô hình 3D trong robot học.
    
    \item \textbf{Chương 5: Khung phương pháp đề xuất Safe-GS (The Proposed Framework):} Trình bày chi tiết toàn bộ đóng góp kỹ thuật của đề tài: từ bước nhận diện vật thể bằng câu lệnh, tối ưu hóa Gaussian ép sát khung hình học, cơ chế quan sát góc khuất bằng Entropy, cho đến bước bù khuyết toàn diện 360$^\circ$, làm sạch mặt bàn và xuất mô hình sang môi trường mô phỏng MuJoCo.
    
    \item \textbf{Chương 6: Kết luận và Hướng phát triển (Conclusion and Future Work):} Tổng kết các đóng góp chính của đề tài, thẳng thắn nhìn nhận các hạn chế kỹ thuật còn tồn đọng và đề xuất các hướng mở rộng tiếp theo cho các bài toán tương tác robot trong tương lai.
\end{itemize}
